% GENERATED FILE. Edit canonical lessons, then run npm run generate:books.
% canonical-source-hash: fnv1a64:b7b7cbd20b7a59bf
% canonical-lessons: MR-C393-sabab, MR-C393-bahana, MR-C393-dhudkavne, MR-C393-vinanti-karne, MR-C393-pathpurava-karne

\chapter{Excuses and Requests}
\label{ch:mr-excuses-and-requests}
\hlchaptermodality{393}

\begin{chapteropening}
\textbf{By the end of this chapter:} \emph{I can talk about excuses, pretexts and requests.}

Complete the last lesson of chapter 393: I can talk about excuses, pretexts and requests.
\end{chapteropening}

\section[\texorpdfstring{\mr{सबब} (sabab)}{सबब (sabab)}]{\mr{सबब} (sabab) — an excuse}

\label{lesson:MR-C393-sabab}



\begin{quote}
Before the new one: say the Marathi for to follow, then the Marathi for to embrace.
\end{quote}

\subsection*{You'll want to know: \mr{सबब}}
\textbf{\mr{सबब}} — \emph{sabab} — \textquotedblleft{}an excuse\textquotedblright{}.

\begin{quote}
\mr{उशिरा} \mr{येण्याची} \mr{त्याची} \mr{सबब} \mr{कोणालाच} \mr{पटली} \mr{नाही}.
\end{quote}

\begin{quote}
\emph{uśirā yeṇyācī tyācī sabab koṇālāc paṭlī nāhī.}
\end{quote}

\begin{quote}
Nobody believed his excuse for coming late.
\end{quote}

Feminine.

Feminine: \mr{सबब} \mr{सांगणे}, to make an excuse.

\begin{cousinweb}
\textbf{\mr{सबब}} (\emph{sabab}) is a loan from Persian, usually traced to Arabic \emph{sabab}, \textquotedblleft{}cause, reason\textquotedblright{}.
\end{cousinweb}

\begin{grammarlens}[title={What you've built}]
One more word for saying what happens if.
\end{grammarlens}

\subsection*{Your turn}
\begin{itemize}
\raggedright
  \item \emph{Say it:} \emph{sabab}
  \item \emph{Say it:} \emph{sabab}, once more
  \item \emph{Say it:} \emph{uśirā yeṇyācī tyācī sabab koṇālāc paṭlī nāhī.} — Nobody believed his excuse for coming late.
  \item \emph{Recall:} say the Marathi for to follow, then the Marathi for to embrace, then say \emph{sabab} again
\end{itemize}

\begin{tcolorbox}[breakable,colback=teal!4,colframe=teal!35!black,title={Before you move on}]
What does \mr{सबब} mean? (\textquotedblleft{}An excuse\textquotedblright{}.) Say it once more.
\end{tcolorbox}
\section[\texorpdfstring{\mr{बहाणा} (bahāṇā)}{बहाणा (bahāṇā)}]{\mr{बहाणा} (bahāṇā) — a pretext}

\label{lesson:MR-C393-bahana}



\begin{quote}
Before the new one: say the Marathi for to embrace, then the Marathi for an excuse.
\end{quote}

\subsection*{You'll want to know: \mr{बहाणा}}
\textbf{\mr{बहाणा}} — \emph{bahāṇā} — \textquotedblleft{}a pretext\textquotedblright{}.

\begin{quote}
\mr{पोटदुखीचा} \mr{बहाणा} \mr{करून} \mr{तो} \mr{शाळेत} \mr{गेला} \mr{नाही}.
\end{quote}

\begin{quote}
\emph{poṭdukhīcā bahāṇā karūn to śāḷet gelā nāhī.}
\end{quote}

\begin{quote}
Pretending to have a stomach ache, he did not go to school.
\end{quote}

Masculine.

X \mr{चा} \mr{बहाणा} \mr{करणे}, to pretend X.

\begin{cousinweb}
\textbf{\mr{बहाणा}} (\emph{bahāṇā}) is a loan from Persian \emph{bahāna}, \textquotedblleft{}pretext, excuse\textquotedblright{}.
\end{cousinweb}

\begin{grammarlens}[title={What you've built}]
One more word for saying what happens if.
\end{grammarlens}

\subsection*{Your turn}
\begin{itemize}
\raggedright
  \item \emph{Say it:} \emph{bahāṇā}
  \item \emph{Say it:} \emph{bahāṇā}, once more
  \item \emph{Say it:} \emph{poṭdukhīcā bahāṇā karūn to śāḷet gelā nāhī.} — Pretending to have a stomach ache, he did not go to school.
  \item \emph{Recall:} say the Marathi for to embrace, then the Marathi for an excuse, then say \emph{bahāṇā} again
\end{itemize}

\begin{tcolorbox}[breakable,colback=teal!4,colframe=teal!35!black,title={Before you move on}]
What does \mr{बहाणा} mean? (\textquotedblleft{}A pretext\textquotedblright{}.) Say it once more.
\end{tcolorbox}
\section[\texorpdfstring{\mr{धुडकावणे} (dhuḍkāvṇe)}{धुडकावणे (dhuḍkāvṇe)}]{\mr{धुडकावणे} (dhuḍkāvṇe) — to brush aside; to spurn}

\label{lesson:MR-C393-dhudkavne}



\begin{quote}
Before the new one: say the Marathi for an excuse, then the Marathi for a pretext.
\end{quote}

\subsection*{You'll want to know: \mr{धुडकावणे}}
\textbf{\mr{धुडकावणे}} — \emph{dhuḍkāvṇe} — \textquotedblleft{}to brush aside; to spurn\textquotedblright{}.

\begin{quote}
\mr{त्याने} \mr{माझा} \mr{सल्ला} \mr{धुडकावला}.
\end{quote}

\begin{quote}
\emph{tyāne mājhā sallā dhuḍkāvlā.}
\end{quote}

\begin{quote}
He brushed aside my advice.
\end{quote}

Past, with -\mr{ने} (he …): \mr{त्याने} \mr{सल्ला} \mr{धुडकावला}, \emph{tyāne sallā dhuḍkāvlā}.

Stronger than \mr{नाकारणे}.

\begin{cousinweb}
\textbf{\mr{धुडकावणे}} (\emph{dhuḍkāvṇe}) is a Marathi verb whose origin is not settled.
\end{cousinweb}

\begin{grammarlens}[title={What you've built}]
One more word for saying what happens if.
\end{grammarlens}

\subsection*{Your turn}
\begin{itemize}
\raggedright
  \item \emph{Say it:} \emph{dhuḍkāvṇe}
  \item \emph{Say it:} \emph{dhuḍkāvṇe}, once more
  \item \emph{Say it:} \emph{tyāne mājhā sallā dhuḍkāvlā.} — He brushed aside my advice.
  \item \emph{Recall:} say the Marathi for an excuse, then the Marathi for a pretext, then say \emph{dhuḍkāvṇe} again
\end{itemize}

\begin{tcolorbox}[breakable,colback=teal!4,colframe=teal!35!black,title={Before you move on}]
What does \mr{धुडकावणे} mean? (\textquotedblleft{}To brush aside; to spurn\textquotedblright{}.) Say it once more.
\end{tcolorbox}
\section[\texorpdfstring{\mr{विनंती} \mr{करणे} (vinantī karṇe)}{विनंती करणे (vinantī karṇe)}]{\mr{विनंती} \mr{करणे} (vinantī karṇe) — to request}

\label{lesson:MR-C393-vinanti-karne}



\begin{quote}
Before the new one: say the Marathi for a pretext, then the Marathi for to brush aside.
\end{quote}

\subsection*{You'll want to know: \mr{विनंती} \mr{करणे}}
\textbf{\mr{विनंती} \mr{करणे}} — \emph{vinantī karṇe} — \textquotedblleft{}to request\textquotedblright{}.

\begin{quote}
\mr{आम्ही} \mr{त्यांना} \mr{थोडं} \mr{थांबण्याची} \mr{विनंती} \mr{केली}.
\end{quote}

\begin{quote}
\emph{āmhī tyānnā thoḍaṁ thāmbaṇyācī vinantī kelī.}
\end{quote}

\begin{quote}
We asked them to wait a little.
\end{quote}

Past, with -\mr{ने} (he …): \mr{त्याने} \mr{विनंती} \mr{केली}, \emph{tyāne vinantī kelī}.

The person asked takes -\mr{ला}, the thing asked for -\mr{ची}.

\begin{cousinweb}
\textbf{\mr{विनंती} \mr{करणे}} (\emph{vinantī karṇe}) joins \mr{विनंती}, \textquotedblleft{}a request\textquotedblright{}, a Marathi noun whose origin is not settled, to \mr{करणे}, \textquotedblleft{}to do\textquotedblright{}.
\end{cousinweb}

\begin{grammarlens}[title={What you've built}]
One more word for saying what happens if.
\end{grammarlens}

\subsection*{Your turn}
\begin{itemize}
\raggedright
  \item \emph{Say it:} \emph{vinantī karṇe}
  \item \emph{Say it:} \emph{vinantī karṇe}, once more
  \item \emph{Say it:} \emph{āmhī tyānnā thoḍaṁ thāmbaṇyācī vinantī kelī.} — We asked them to wait a little.
  \item \emph{Recall:} say the Marathi for a pretext, then the Marathi for to brush aside, then say \emph{vinantī karṇe} again
\end{itemize}

\begin{tcolorbox}[breakable,colback=teal!4,colframe=teal!35!black,title={Before you move on}]
What does \mr{विनंती} \mr{करणे} mean? (\textquotedblleft{}To request\textquotedblright{}.) Say it once more.
\end{tcolorbox}
\section[\texorpdfstring{\mr{पाठपुरावा} \mr{करणे} (pāṭhpurāvā karṇe)}{पाठपुरावा करणे (pāṭhpurāvā karṇe)}]{\mr{पाठपुरावा} \mr{करणे} (pāṭhpurāvā karṇe) — to follow up; to pursue}

\label{lesson:MR-C393-pathpurava-karne}



\begin{quote}
Before the new one: say the Marathi for to brush aside, then the Marathi for to request.
\end{quote}

\subsection*{You'll want to know: \mr{पाठपुरावा} \mr{करणे}}
\textbf{\mr{पाठपुरावा} \mr{करणे}} — \emph{pāṭhpurāvā karṇe} — \textquotedblleft{}to follow up; to pursue\textquotedblright{}.

\begin{quote}
\mr{तक्रारीचा} \mr{पाठपुरावा} \mr{केल्यावर} \mr{काम} \mr{झालं}.
\end{quote}

\begin{quote}
\emph{takrārīcā pāṭhpurāvā kelyāvar kām jhālaṁ.}
\end{quote}

\begin{quote}
After the complaint was followed up, the work got done.
\end{quote}

Past, with -\mr{ने} (he …): \mr{त्याने} \mr{पाठपुरावा} \mr{केला}, \emph{tyāne pāṭhpurāvā kelā}.

What is pursued takes -\mr{चा}: \mr{तक्रारीचा} \mr{पाठपुरावा}.

\begin{cousinweb}
\textbf{\mr{पाठपुरावा} \mr{करणे}} (\emph{pāṭhpurāvā karṇe}) joins \mr{पाठपुरावा}, \textquotedblleft{}pursuit, follow-up\textquotedblright{}, a Marathi noun whose origin is not settled, to \mr{करणे}, \textquotedblleft{}to do\textquotedblright{}.
\end{cousinweb}

\begin{grammarlens}[title={What you've built}]
One more word for saying what happens if.
\end{grammarlens}

\subsection*{Your turn}
\begin{itemize}
\raggedright
  \item \emph{Say it:} \emph{pāṭhpurāvā karṇe}
  \item \emph{Say it:} \emph{pāṭhpurāvā karṇe}, once more
  \item \emph{Say it:} \emph{takrārīcā pāṭhpurāvā kelyāvar kām jhālaṁ.} — After the complaint was followed up, the work got done.
  \item \emph{Recall:} say the Marathi for to brush aside, then the Marathi for to request, then say \emph{pāṭhpurāvā karṇe} again
\end{itemize}

\begin{tcolorbox}[breakable,colback=teal!4,colframe=teal!35!black,title={Before you move on}]
What does \mr{पाठपुरावा} \mr{करणे} mean? (\textquotedblleft{}To follow up; to pursue\textquotedblright{}.) Say it once more.
\end{tcolorbox}
